\renewcommand{\author}{張晴昀 (B10902064)\\黃允謙 (B10902067)}
\renewcommand{\title}{Project 2: Image Stitching}
\newcommand{\course}{Digital Visual Effects (NTU CSIE, Spring 2025)}

\renewcommand\tabularxcolumn[1]{>{\Centering}m{#1}}

\begin{document}
\maketitle


\section{Description}

This project aims to stitch a set of images into a seamless panorama. The 
process involves several key image processing steps, including feature detection, 
description, and matching. Once corresponding features are matched across images, 
further operations such as warping, blending, alignment and cropping are 
performed to generate the final panoramic output. Image stitching enhances both 
the effective resolution and the field of view, making it a powerful technique 
for producing wide-angle visual representations.


\section{Processing Flow}
\begin{enumerate}
  \item Feature Detection
  \begin{itemize}
    \item Moravec Corner Detector
    \item Harris Corner Detector (Bonus)
  \end{itemize}
  \item Feature Descriptor
  \item Feature Matching
  \item Image matching
  \item Blending
  \item End-to-end Alignment (Bonus)
  \item Rectangling (Bonus)
\end{enumerate}


\section{Feature Detection}
\subsection{Moravec Corner Detector}

We first implement the simple Moravec corner detector. For each pixel, the 
algorithm calculates the intensity difference between its patch and a shifted 
version. Then, it uses the minimum difference among four shifted directions as 
the pixel's energy. A corner exhibits drastic intensity changes in all 
directions, and thus results in higher energy.
$$
  E(\mathbf{u}) = 
  \sum_{X} w(\mathbf{x}) [I(\mathbf x + \mathbf u) - I(\mathbf x)]^2
$$

\subsection{Harris Corner Detector}

However, Moravec has several problems. First, it uses a binary window function 
for the patch, which makes it more susceptible to noise. Second, it only 
considers shifts in eight directions, potentially missing corners at other angles.

The Harris corner detector introduces an alternative method for computing the 
energy to address these issues. It uses a Gaussian function to produce a 
smoother window for the patch. Then, it applies Taylor expansion to approximate 
the intensity difference for a given shift direction, which can be expressed as:
$$
  E(\mathbf{u}) \approxeq \mathbf{u}^\intercal\mathbf{M}\mathbf{u}, \quad 
  \mathbf{M} = \sum_{X} w(\mathbf x) \begin{bmatrix}
    I_x^2 & I_xI_y \\ I_y^2 & I_y^2
  \end{bmatrix}
$$
Since $E$ is a quadratic function of the shift vector $\mathbf{u}$, we can 
analyze its shape by examining the properties of $\mathbf{M}$. As a result, 
Harris considered the eigenvalues of $\mathbf{M}$. When both eigenvalues are 
large, it indicates that $E$ will be large in all directions $\mathbf{u}$, 
suggesting the presence of a corner. Harris proposed the following energy function:
$$  
  E = \mathrm{det}\mathbf{M} - k (\mathrm{tr}\mathbf{M})^2
  = \lambda_1\lambda_2 - k(\lambda_1 + \lambda_2)^2
$$

\subsection{Comparison}

In Figure \ref{fig:feature-detection-comparison}, we demonstrate the results of 
the Moravec and Harris algorithms. In all the images, redder colors indicate 
higher values, while bluer colors indicate lower values.

Both algorithms performs well in identifying the corners or possible feature 
locations. In the Harris algorithm column, we observe that the horizontal and 
vertical lines appear in blue, and the corners appear in red, which aligns with 
our expectations. However, we did not anticipate that a diagonal line would 
appear in red and be classified as a corner. 

\begin{figure}[htbp]
    \centering

    \begin{subfigure}[b]{0.3\textwidth}
        \centering\textbf{Image}
    \end{subfigure}
    \hfill
    \begin{subfigure}[b]{0.3\textwidth}
        \centering\textbf{Moravec}
    \end{subfigure}
    \hfill
    \begin{subfigure}[b]{0.3\textwidth}
        \centering\textbf{Harris}
    \end{subfigure}
    \vskip1ex

    \begin{subfigure}[b]{0.3\textwidth}
        \centering\includegraphics[width=\textwidth]{geometry.jpg}
    \end{subfigure}
    \hfill
    \begin{subfigure}[b]{0.3\textwidth}
        \centering\includegraphics[width=\textwidth]{geometry-moravec.jpg}
    \end{subfigure}
    \hfill
    \begin{subfigure}[b]{0.3\textwidth}
        \centering\includegraphics[width=\textwidth]{geometry-harris.jpg}
    \end{subfigure}
    \vskip\baselineskip

    \begin{subfigure}[b]{0.3\textwidth}
        \centering\includegraphics[width=\textwidth]{toyosato.jpg}
    \end{subfigure}
    \hfill
    \begin{subfigure}[b]{0.3\textwidth}
        \centering\includegraphics[width=\textwidth]{toyosato-moravec.jpg}
    \end{subfigure}
    \hfill
    \begin{subfigure}[b]{0.3\textwidth}
        \centering\includegraphics[width=\textwidth]{toyosato-harris.jpg}
    \end{subfigure}
    \vskip\baselineskip

    \begin{subfigure}[b]{0.3\textwidth}
        \centering\includegraphics[width=\textwidth]{nagoya.jpg}
    \end{subfigure}
    \hfill
    \begin{subfigure}[b]{0.3\textwidth}
        \centering\includegraphics[width=\textwidth]{nagoya-moravec.jpg}
    \end{subfigure}
    \hfill
    \begin{subfigure}[b]{0.3\textwidth}
        \centering\includegraphics[width=\textwidth]{nagoya-harris.jpg}
    \end{subfigure}
    \caption{\small Comparison of Moravec and Harris Algorithms}
    \label{fig:feature-detection-comparison}
\end{figure}

\subsection{Threshold}

To find the thresholds of the energy functions, we attempted the following two 
methods.

\subsubsection{Otsu's Method}
When finding a threshold in a one-dimensional distribution, it is common to 
apply Otsu's method. This algorithm separates values into two classes, aiming to 
minimize the intra-class variances. Otsu's method performs optimally when the 
values follow a bimodal distribution.

In our application of Otsu's method, it may return an acceptable threshold. 
However, as shown in Figure \ref{fig:feature-value-distribution}, we observe 
that the values from both algorithms do not follow a bimodal distribution. This
makes the use of Otsu's method undesirable in this case.

\begin{figure}[htbp]
    \centering
    \begin{subfigure}[b]{0.49\textwidth}
        \centering\includegraphics[width=\textwidth]{moravec-values.png}
    \end{subfigure}
    \hfill
    \begin{subfigure}[b]{0.49\textwidth}
        \centering\includegraphics[width=\textwidth]{harris-values.png}
    \end{subfigure}
    \caption{\small
        The energy value distributions of Moravec and Harris algorithms on the 
        image of Nagoya Castle. The red dotted line indicates Otsu's threshold.
    }
    \label{fig:feature-value-distribution}
\end{figure}

\subsubsection{Heuristic}
We developed another heuristic method to determine the threshold. As shown in 
Figure \ref{fig:feature-value-distribution}, both distributions exhibit a spike.
We believe this is the flat areas in the image, which should be a major pattern 
in any image. When we include values above this spike, we obtain a moderate 
number of features. However, if we include the spike itself, nearly the entire 
image becomes a set of feature candidates. Note that the y-axis is on a 
logarithmic scale, so the spike is more than ten times larger than the other 
value counts. 

We decided to use this spike value as the threshold, and it turned out to be 
quite effective. The algorithm proceeds as follows: first, we quantize the 
values into 256 bins. Then, we identify the bin with the highest count of values 
and set it as the threshold. 

As shown in Figure \ref{fig:feature-value-distribution}, this threshold is 
slightly smaller then the one determined by Otsu's method. It includes more 
possible candidates while avoiding the includsion of the entire image. As shown 
in Figure \ref{fig:feature-thresholds}, Otsu's method fails to detect potential 
features on the bottom-left roof tiles, while our heuristic method successfully 
captures them.

\begin{figure}[htbp]
    \centering
    \begin{subfigure}[b]{0.32\textwidth}
        \centering\includegraphics[width=\textwidth]{features-otsu.jpg}
        \caption{Otsu's Method}
    \end{subfigure}
    \hfill
    \begin{subfigure}[b]{0.32\textwidth}
        \centering\includegraphics[width=\textwidth]{features-heuristic.jpg}
        \caption{Heuristic}
    \end{subfigure}
    \hfill
    \begin{subfigure}[b]{0.32\textwidth}
        \centering
\includegraphics[width=\textwidth]{features-heuristic-predecessor.jpg}
        \caption{Heuristic include spike}
    \end{subfigure}
    \caption{\small
        Features of the three different threshold methods. The Heuristic method 
        finds the most possible candidates while avoiding including the entire 
        image. 
    }
    \label{fig:feature-thresholds}
\end{figure}

\subsection{Non-maximum Suppression}

Since we might have an small area of pixels all exceeding the threshold, making 
all of them feature candidates. However, having many overlapping features is 
undesirable. Thus we include non-maximum suppression. The idea is to select the 
pixel with the maximum value in a small patch as the only feature. 

In practice, we apply a maximum filter mask, that only include pixels which 
appears to be the maximum value within a given distance. We then use this mask 
to filter feature candidates to produce a sparse result. As shown in Figure 
\ref{fig:feature-filter}, the unfiltered version marks nearly all roof tiles as 
features. After applying non-maximum suppression, the features become more 
spatially separated and tend to correspond to corner points in the tile pattern.

\begin{figure}[htbp]
    \centering
    \begin{subfigure}[b]{0.48\textwidth}
        \centering\includegraphics[width=\textwidth]{features-unfiltered.jpg}
        \caption{Unfiltered}
    \end{subfigure}
    \hfill
    \begin{subfigure}[b]{0.48\textwidth}
        \centering\includegraphics[width=\textwidth]{features-filtered.jpg}
        \caption{Filtered}
    \end{subfigure}
    \caption{\small Filtering effect of applying non-maximum suppression.}
    \label{fig:feature-filter}
\end{figure}

\subsection{Interactive Tool}

To better visualize and test the effect of threshold adjustment. We developed
a small interactive tool that displays real-time feature detection results as 
the user adjusts the threshold using a slider. With this tool, we observed that 
in many images, there is a significant jump in the number of detected features 
at a certain threshold. This observation prompted us to examine the distribution 
of the energy values more closely, ultimately leading to the development of our 
heuristic thresholding method. 

\begin{figure}[htbp]
    \centering 
    \includegraphics[width=0.7\textwidth]{interactive-tool.png}
    \caption{Interactive tool}
    \label{fig:interactive-tool}
\end{figure}

\section{Feature Matching}

\subsection{Feature Descriptor}

Since our photos were taken using a tripod and identical camera settings, there 
is minimal scaling or rotation between images. As a result, we found that simply 
using a small patch around each feature point works well for this project.

We experimented with patch sizes $l \in \set{3, 5, 7}$ across two types of 
images -- BGR color channels and grayscale -- resulting in a total of six 
combinations. Each patch is flattened into a one-dimensional vector to serve as 
the feature descriptor. However, using $l = 7$ with BGR color channels results 
in a descriptor dimension of $7\times 7\times 3 = 147$, which exceeds the 
memory capacity of our device during the acceleration step, causing an
out-of-memory error. Therefore, we were unable to obtain results for this 
combination. 

\subsection{Matching}

To establish correspondences between two images, we first detect feature points 
in the first image and extract a descriptor for each feature. Then, we compute 
descriptors for every pixel in the second image and attempt to find the best 
match for each feature. We define the best match as the pair of descriptors with 
the minimum distance.

Assume the first image is on the left and the second image be on the right. 
Since the leftmost region of the first image and the rightmost region of the 
second image do not overlap, including these areas in the matching process will 
introduces unnecessary computation and disturbance. To address this, we 
introduce a parameter that allows the user to specify the overlapping ratio 
between the two images. This not only reduces the search space and speeds up the 
matching process, but also decreases the number of unmatched features, thereby 
improving overall matching accuracy.

\subsubsection{Shift Heatmap}

It is difficult to determine if all matches are correct without manually 
checking each one. Therefore, we validate the matches by checking if they are 
consistent with each other. If all the matches are correct, the shifts in the 
matched coordinates should exhibit similar patterns due to the camera rotation. 
To visualize this, we plot all the shifts on a heatmap to determine if they are 
centralized around a single point. Conversely, a set of matches is considered 
incorrect if the shifts are scattered. We use this visualization technique to 
evaluate and compare the following results.

The $x$ and $y$ axes of the heatmap represent the coordinate shifts of the 
matched feature points. To better visualize the results, we zoomed into the 
heatmap. Additionally, the shift count at each pixel location is passed through 
the following transformation function $f$, which emphasizes lower values and 
enhances the visibility of scattered patterns:
$$
    f(x) = \log(1+100x)
$$

\subsubsection{Acceleration}

To accelerate the matching process, we employed search data structures such as 
the K-D Tree. We spent additional time constructing a K-D tree for all 
descriptors in the second image, which significantly reduced the time required 
to find the nearest vector. We also experimented with Ball Tree and discovered 
that K-D Tree performed slightly better in our case. 

\begin{table}[htbp]
    \centering
    \begin{minipage}{0.48\textwidth}
        \centering
        \begin{tabular}{|c|c|c|}
            \hline
            \textbf{Patch} & \textbf{BGR Color} & \textbf{Grayscale} \\
            \hline
            3 & 51 + 0 = 51 & 27 + 0 = 27 \\
            \hline
            5 & 127 + 17 = 144 & 62 + 6 = 68 \\
            \hline
            7 & \textit{out-of-memory} & 111 + 70 = 181 \\
            \hline
        \end{tabular}
        \caption*{K-D Tree}
    \end{minipage}
    \hfill 
    \begin{minipage}{0.48\textwidth}
        \centering
        \begin{tabular}{|c|c|c|}
            \hline
            \textbf{Patch} & \textbf{BGR Color} & \textbf{Grayscale} \\
            \hline
            3 & 45 + 5 = 50 & 25 + 2 = 26 \\
            \hline
            5 & 117 + 62 = 180 & 57 + 24 = 81 \\
            \hline
            7 & \textit{out-of-memory} & 103 + 106 = 209 \\
            \hline
        \end{tabular}
        \caption*{Ball Tree}
    \end{minipage}
    \label{tab:match-time}
    \caption{\small Construction and querying times for 500 features across 
    different descriptor combinations, measured in seconds.}
\end{table}

\subsubsection{Match Filtering} \label{match-filtering}
Sometimes, no correct match can be found for a feature in the first image. This 
may be due to object movement between shots or simply because the feature is not 
visible in the second image. In such cases, selecting the nearest but incorrect 
match would negatively impact the subsequent image matching process. To address 
this, we apply a unique filter to the matches.

The approach is to accept a match only when there is a clear and unambiguous 
correspondence. If there are two or more similar candidates, we discard the 
feature to avoid mismatches. In practice, we search for the top two nearest 
vectors in the second image’s descriptors and accept the match if the ratio of 
their distances is smaller than a given threshold. 

As shown in Table \ref{tab:feature-match-threshold}, we observe a significant amount of 
mismatches when the filter is not applied (i.e., threshold = 1.0). The 
corresponding heatmap also displays many scattered points, indicating 
inconsistency among the shift vectors. After experimenting with various 
parameters, we found that setting the threshold to 0.8 provides the best balance 
between obtaining a sufficient number of matched features and minimizing 
mismatches.

$$
    \frac{\text{nearest distance}}{\text{second-nearest distance}} < 
    \text{threshold} = 0.8
$$

\begin{table}[htbp]
    \centering
    \renewcommand{\arraystretch}{1.5}
    \begin{tabular}{|c|c|>{\centering\arraybackslash}m{0.34\textwidth}|>{\centering\arraybackslash}m{0.34\textwidth}|}
        \hline
        \textbf{Thresh} & \textbf{\#Match} & \textbf{Matchings} & \textbf{Heatmap} \\
        \hline
        1.0 & 500 
            & \includegraphics[width=0.34\textwidth]{match-filter/thresh-10-match.jpg}
            & \includegraphics[width=0.34\textwidth]{match-filter/thresh-10-heatmap.png} \\
        \hline
        0.9 & 271 
            & \includegraphics[width=0.34\textwidth]{match-filter/thresh-09-match.jpg} 
            & \includegraphics[width=0.34\textwidth]{match-filter/thresh-09-heatmap.png} \\
        \hline
        0.8 & 179
            & \includegraphics[width=0.34\textwidth]{match-filter/thresh-08-match.jpg} 
            & \includegraphics[width=0.34\textwidth]{match-filter/thresh-08-heatmap.png} \\
        \hline
        0.7 & 118
            & \includegraphics[width=0.34\textwidth]{match-filter/thresh-07-match.jpg} 
            & \includegraphics[width=0.34\textwidth]{match-filter/thresh-07-heatmap.png} \\
        \hline
    \end{tabular}
    \caption{Comparison of match filtering and heatmaps across varying threshold values.}
    \label{tab:feature-match-threshold}
\end{table}

\subsection{Results}

We compared the results of the five possible feature descriptor combinations, 
excluding one that could not be computed due to an out-of-memory error. All 
matching results discussed below are filtered using a threshold of 0.8 as 
mentioned in section \ref{match-filtering}.

As shown in Table \ref{tab:descriptor-match-count}, using a patch size of 3 or 5 
does not significantly affect the number of matched features. However, 
increasing the patch size to 7 nearly doubles the match count. We hypothesize 
that a higher-dimensional descriptor provides more distinctive information, 
increasing the likelihood of finding a unique match that passes the threshold 
filter. Nonetheless, the BGR color descriptor with size 5 has a higher 
dimensionality (75) than the grayscale descriptor with size 7 (49). This 
suggests that the three color channels in the BGR space may be highly 
correlated, reducing the effective dimensionality of the descriptor.

In addition to quantity, the quality of the matches must also be considered. As 
shown in Tables \ref{tab:desciptor-matchings} and \ref{tab:desciptor-heatmaps}, 
a patch size of 3 produces unstable results, with many visible mismatches and 
highly scattered shift heatmaps. Increasing the patch size to 5 improves the 
performance, with BGR descriptors outperforming grayscale. However, the 
grayscale descriptor with patch size 7 clearly yields the most consistent 
matches and a well-centered heatmap, indicating both high match quantity and 
quality. Therefore, we choose this combination as our final descriptor.

\begin{table}[htbp]
    \centering
    \begin{tabular}{|c|c|c|}
        \hline
        \textbf{Patch} & \textbf{BGR Color} & \textbf{Grayscale} \\
        \hline
        3 & 74 & 86 \\
        \hline
        5 & 96 & 101 \\
        \hline
        7 & \textit{OOM} & 179 \\
        \hline
    \end{tabular}
    \caption{Number of matched features of different descriptor combinations.}
    \label{tab:descriptor-match-count}
\end{table}

\begin{table}[htbp]
    \centering
    \renewcommand{\arraystretch}{1.5}
    \begin{tabularx}{0.95\textwidth}{|c|X|X|}
        \hline
        \textbf{Patch} & \textbf{BGR Color} & \textbf{Grayscale} \\
        \hline
        3 & \includegraphics[width=0.9\linewidth]{descriptor/color-3-match.jpg} 
          & \includegraphics[width=0.9\linewidth]{descriptor/gray-3-match.jpg} \\
        \hline
        5 & \includegraphics[width=0.9\linewidth]{descriptor/color-5-match.jpg} 
          & \includegraphics[width=0.9\linewidth]{descriptor/gray-5-match.jpg} \\
        \hline
        7 & \textit{out-of-memory}
          & \includegraphics[width=0.9\linewidth]{descriptor/gray-7-match.jpg} \\
        \hline
    \end{tabularx}
    \caption{Matching results of different descriptors combinations.}
    \label{tab:desciptor-matchings}
\end{table}

\begin{table}[htbp]
    \centering
    \renewcommand{\arraystretch}{1.5}
    \begin{tabularx}{0.95\textwidth}{|c|X|X|}
        \hline
        \textbf{Patch} & \textbf{BGR Color} & \textbf{Grayscale} \\
        \hline
        3 & \includegraphics[width=0.9\linewidth]{descriptor/color-3-heatmap.png} 
          & \includegraphics[width=0.9\linewidth]{descriptor/gray-3-heatmap.png} \\
        \hline
        5 & \includegraphics[width=0.9\linewidth]{descriptor/color-5-heatmap.png} 
          & \includegraphics[width=0.9\linewidth]{descriptor/gray-5-heatmap.png} \\
        \hline
        7 & \textit{out-of-memory}
          & \includegraphics[width=0.9\linewidth]{descriptor/gray-7-heatmap.png} \\
        \hline
    \end{tabularx}
    \caption{Shift heatmaps of different descriptors combinations.}
    \label{tab:desciptor-heatmaps}
\end{table}

\section{Image Matching}

To perform image matching, we estimate a 2D translation transformation that aligns 
two images based on a set of matched feature points using the RANSAC algorithm. 
The translation model assumes that the transformation between corresponding points 
in the two images can be described by a vector $(d_x, d_y)$, where each point 
$(x, y)$ in the first image maps to $(x + d_x, y + d_y)$ in the second image. 
RANSAC is used to estimate this translation by iteratively selecting minimal 
subsets of point correspondences (in this case, one pair is sufficient for 
translation), computing the translation vector from the subset, and 
identifying inliers that agree with the estimated model within a predefined 
threshold. The model with the highest number of inliers is selected as the final 
translation estimate.

\subsection{Threshold}

In Figure~\ref{fig:image-matching-comparison}, we apply different thresholds 
$(5, 10, 50, 100)$ in the RANSAC algorithm for comparison. The feature points 
used for matching are extracted using the Harris Corner Detector. For the Day 
panoramas, when the threshold is set to 100, neither building aligns well. At 
a threshold of 50, the left building aligns well, but the right one remains 
misaligned. When the threshold is reduced to 10, both buildings are well aligned, 
resulting in a coherent panorama. Further lowering the threshold to 5 produces a 
result that is visually similar to the one obtained with a threshold of 10. 
For the Night panoramas, the alignment appears generally consistent across all 
thresholds. However, a slight shift is observed in the roof region of each 
panorama, and this shift becomes less noticeable as the threshold decreases.
The Night panorama with threshold of 100 shows better alignment than the 
corresponding Day panorama. This difference may be attributed to the presence 
of moving people in the Day images. Some feature points may have been detected 
on those people, which do not occupy consistent positions across the images.

\begin{figure}[htbp]
  \centering

  \begin{subfigure}[b]{\textwidth}
    \centering
    \begin{minipage}{0.45\textwidth}
      \centering\textbf{Day}
    \end{minipage}
    \hfil
    \begin{minipage}{0.45\textwidth}
      \centering\textbf{Night}
    \end{minipage}
  \end{subfigure}
  \vskip1ex

  \begin{subfigure}[b]{\textwidth}
    \centering
    \begin{minipage}{0.45\textwidth}
      \includegraphics[width=\textwidth]{day_100.pdf}
    \end{minipage}
    \hfil
    \begin{minipage}{0.45\textwidth}
      \includegraphics[width=\textwidth]{night_100.pdf}
    \end{minipage}
    \caption{threshold = 100}
  \end{subfigure}
  \vskip\baselineskip

  \begin{subfigure}[b]{\textwidth}
    \centering
    \begin{minipage}{0.45\textwidth}
      \includegraphics[width=\textwidth]{day_50.pdf}
    \end{minipage}
    \hfil
    \begin{minipage}{0.45\textwidth}
      \includegraphics[width=\textwidth]{night_50.pdf}
    \end{minipage}
    \caption{threshold = 50}
  \end{subfigure}
  \vskip\baselineskip

  \begin{subfigure}[b]{\textwidth}
    \centering
    \begin{minipage}{0.45\textwidth}
      \includegraphics[width=\textwidth]{day_10.pdf}
    \end{minipage}
    \hfil
    \begin{minipage}{0.45\textwidth}
      \includegraphics[width=\textwidth]{night_10.pdf}
    \end{minipage}
    \caption{threshold = 10}
  \end{subfigure}
  \vskip\baselineskip
  
  \begin{subfigure}[b]{\textwidth}
    \centering
    \begin{minipage}{0.45\textwidth}
      \includegraphics[width=\textwidth]{day_5.pdf}
    \end{minipage}
    \hfil
    \begin{minipage}{0.45\textwidth}
      \includegraphics[width=\textwidth]{night_5.pdf}
    \end{minipage}
    \caption{threshold = 5}
  \end{subfigure}

  \caption{Comparison of RANSAC performance under different thresholds}
  \label{fig:image-matching-comparison}
\end{figure}

\section{Blending}

To seamlessly combine the aligned images, we apply linear blending in the 
overlapping region. After estimating the translation, the overlapping area between 
two images is computed. In this region, pixel values are blended using a weighted 
sum, where the weights vary linearly across the overlap to ensure a smooth 
transition. Specificallym for each pixel $p$ in the region, the final value is 
computed as 
$$
I(p) = \alpha(p_1) \cdot I_1(p_1) + (1 - \alpha(p_1)) \cdot I_2(p_2),
$$ 
where $I_1$ and $I_2$ are the intensity values from the two images, $p_1$ and $p_2$ 
are the corresponding pixels in the two images, and $\alpha(p) \in [0, 1]$ is a 
spatially varying weight based on the pixel's relative position within the 
overlapping area. In our implementation, we asssume the first image is on the left 
and the second image is on the right. The weight $\alpha(p_1)$ is determined based 
on the pixel's horizontal distance from the right edge of the first image: pixels 
at the edge are assigned a weight of $0$, pixels at the start of the overlapping 
area are assigned $1$, and those in between are weighted linearly between $[0, 1]$.

\section{End-to-end Alignment}

To reduce visual discontinuities caused by vertical shift, we implement an 
end-to-end alignment strategy that distributes the vertical shift across image 
cloumns. After stitching all images to form a panorama, we apply a column-wise  
vertical adjustment. For example, if the resulting panorama is tilted toward the 
lower right, we shift the first column downward by a specified amount, while the 
last column remains unchanged. The vertical shift for columns in between is 
interpolated linearly within the range $[0, \text{shift of the first column}]$. 
This gradual adjustment ensures a smooth vertical transition across the panorama, 
as shown Figure~\ref{fig:end-to-end-alignment}.

\begin{figure}[htbp]
  \centering

  \begin{subfigure}[b]{\textwidth}
    \centering
    \begin{minipage}{0.45\textwidth}
      \centering\textbf{Day}
    \end{minipage}
    \hfil
    \begin{minipage}{0.45\textwidth}
      \centering\textbf{Night}
    \end{minipage}
  \end{subfigure}
  \vskip1ex

  \begin{subfigure}[b]{\textwidth}
    \centering
    \begin{minipage}{0.45\textwidth}
      \includegraphics[width=\textwidth]{day_5.pdf}
    \end{minipage}
    \hfil
    \begin{minipage}{0.45\textwidth}
      \includegraphics[width=\textwidth]{night_5.pdf}
    \end{minipage}
    \caption{Before end-to-end alignment}
  \end{subfigure}
  \vskip\baselineskip

  \begin{subfigure}[b]{\textwidth}
    \centering
    \begin{minipage}{0.45\textwidth}
      \includegraphics[width=\textwidth]{day_5_aligned.pdf}
    \end{minipage}
    \hfil
    \begin{minipage}{0.45\textwidth}
      \includegraphics[width=\textwidth]{night_5_aligned.pdf}
    \end{minipage}
    \caption{After end-to-end alignment}
  \end{subfigure}

  \caption{Effects of the end-to-end alignment}
  \label{fig:end-to-end-alignment}
\end{figure}


\section{Rectangling}

We attempted to implement the algorithm proposed by \textbf{He et al. (2013)}, 
but due to time constraints and the algorithm's complexity, we were only able to 
complete the first stage. This stage involves using seam carving to synthesize 
pixels that expand the panorama into a rectangular shape. The second stage, 
which we did not completed, involves warping the image back and apply a 
carefully designed energy function to guide a smooth reshaping process. 

In this section, we present our findings and results based on implementing only 
the seam carving stage. 

\subsection{Seam Carving}

Seam carving is an algorithm proposed by \textbf{Avidan and Shamir (2007)} that 
primarily focuses on resizing rectangular images. In their original paper, the 
authors define several energy functions to measure the importance of each pixel. 
The algorithm then identifies an 8-connected path from one edge of the image to 
the other that accumulates the least energy, which is referred to as a "seam". 
Depending on whether the user wishes to increase or decrease the image size, the 
seam is either duplicated or removed, respectively. 

The key idea behind seam carving is to modify the less visually noticeable parts
of the image, ensuring the modifications are seamless. One of the energy 
functions proposed by the authors is:
$$
    e(\mathbf I) = \left| \frac{\partial}{\partial x} \mathbf I\right| + 
    \left| \frac{\partial}{\partial y} \mathbf I\right|
$$  
This energy function aims to identify regions with minimal intensity change, 
helping to avoid object edges that are more sensitive to human vision. 

However, in the panorama rectangling problem, the goal is not to insert entire 
row of pixels, but rather fill irregularly shaped gaps. To address this, 
\textbf{He et al.} propose a method where we first identify a line of 
empty pixels along the boundaries of the image. A sub-image is then extracted, 
and the original seam carving technique is applied. This process is repeated 
iteratively: a new empty segment is found, a sub-image is cut, and a seam is 
added, until no more gaps remain, at which point the image is successfully 
filled into a rectangular shape. This iterative process is referred to as 
"local warping."

\subsection{Implementation}

In our implementation, we used the same energy function described earlier. To 
find the seam with the least accumulated energy, we applied the dynamic 
programming technique suggested in the paper, which reduces the time complexity 
to be linear in relation to the image size.

Since segments may appear along all four edges of the image, we need to handle 
four different cases of seam carving, which complicates the implementation. To 
simplify this, we defined an additional \code{Transformation} class. This class 
helps cut the sub-image and transform it such that the segment is always placed 
on the top edge. After applying the seam carving algorithm, we revert the 
transformation to restore the correct coordinates. This approach significantly 
reduces redundant code, making the implementation cleaner and more efficient.

\begin{figure}[htbp]
    \centering\includegraphics[width=0.9\textwidth]{seam-carving.png}
    \caption{An example of seam insertion: the white line represents the 
    duplicated seam, and the pixels above it are shifted upward by one pixel to 
    fill the resulting gap.}
\end{figure}

\subsection{Results}

If we apply only seam carving and local warping, without the subsequent global 
warping, the resulting image will have visible cracks, as shown in Figure 
\ref{fig:day-seam-failed}. This occurs because we apply seam carving only to a 
sub-image rather than the entire image. When we insert a seam, we optimize the 
energy function within the sub-image, but we also create a shift at the 
boundaries between the sub-image and the original image. As we repeatedly insert 
seams into the same sub-image, these shifts accumulate, leading to visible 
artifacts.

However, when we apply this process to a night-time photo, the shifts 
predominantly affect the upper and lower parts of the image, which are dark and 
less visible. As a result, the shifted artifacts are not as noticeable, and the 
final image appears relatively artifact-free, as shown in Figure 
\ref{fig:night-rect}.

\begin{figure}[htbp]
    \centering\includegraphics[width=0.6\textwidth]{day-seam-failed.png}
    \caption{Artifacts produced by seam carving.}
    \label{fig:day-seam-failed}
\end{figure}

\begin{figure}[htbp]
    \centering\includegraphics[width=0.95\textwidth]{night-rect.jpg}
    \caption{Final results after seam carving on the night-time image.}
    \label{fig:night-rect}
\end{figure}

\end{document}